\section{Experimental Design}
\subsection{Dataset and split}
The completed P001 pilot uses Fashion-MNIST, a ten-class collection of 28-by-28 grayscale
fashion images \cite{xiao2017fashion}. The standard source contains 60,000 training and
10,000 test images; P001 intentionally uses a reduced deterministic slice to validate the
research machinery. The frozen split contains 9,000 training, 1,000 validation, and 2,500
test examples. Dataset acquisition, scaling, split IDs, and the post-scaling dataset hash
are recorded. This reduced slice is not presented as publication-scale evidence.

The small slice makes the control path inspectable: a reviewer can reconstruct every split,
rerun every exact replica, and inspect each raw record without requiring a large compute
allocation. It does not establish that the observed variability, deletion behavior, or
method ranking would persist at the scale, augmentation regime, and optimization horizon of
a main benchmark. Fashion-MNIST is therefore used here as a bridge between Tier-0
infrastructure checks and a future registered multi-dataset study, not as the empirical
basis for a general claim about approximate unlearning.

\subsection{Deletion request}
The deletion policy samples one percent of the frozen training split uniformly without
replacement using deletion seed 1009. The resulting immutable manifest contains
\PilotForgetCount{} stable example IDs, leaving 8,910 retained examples. Both lists and
their hashes are stored; the retained list is verified as the exact complement of the
deletion list. This is an IID deletion, not a class-deletion, high-influence, rare-example,
corrupted-label, or memorization stress regime.

IID deletion provides a necessary reference condition because it separates the mechanics of
manifest construction and counterfactual retraining from a selection rule that deliberately
targets unusual training examples. It is not expected to be the most revealing deletion
regime. The experiment schema retains explicit selection interfaces for class-specific,
rare-example, clustered, mislabeled, influence-selected, and sequential requests. Those
regimes should be introduced as registered factors rather than retrofitted after inspecting
the IID pilot, since their selection procedures can change both the difficulty of unlearning
and the appropriate interpretation of forget-set metrics.

\subsection{Model and training}
P001 uses the repository's small convolutional architecture with hidden dimension 64.
Original and exact models are trained for four epochs with AdamW, batch size 128, learning
rate $10^{-3}$, and weight decay $10^{-4}$. Deterministic execution is requested and the
device is selected automatically. The original model uses seed 11; the exact population
uses seeds 2001--2015. All exact models train on the same 8,910-example retained set.

Holding the architecture, preprocessing, training schedule, and retained set fixed makes
initialization and remaining stochastic training effects the intended source of variation
within the exact population. Deterministic settings reduce incidental implementation
variation within a run, but they do not make a single seed representative of all valid
counterfactual trainings. The seed ensemble is therefore a design feature rather than a
replication afterthought: its pairwise distances define the reference scale against which an
approximate output is assessed.

\subsection{Approximate engineering baselines}
The pilot executes fine-tuning, naive gradient ascent, and SCRUB. Fine-tuning and naive
gradient ascent are explicitly simple baselines; they are not described as competitive
leading methods. The archived configuration records one unlearning epoch,
learning rate $2\times10^{-4}$, retain batch size 128, forget batch size 64, retain weight
1.0, forget weight 0.1, and the method-specific controls used by the runner. These choices
validate a common interface and are not the product of a publication-scale tuning study.

The distinction between an engineering baseline and a validated comparison method is
important. Fine-tuning and naive ascent answer whether the evaluation and artifact pipeline
can process simple parameter updates; they do not represent the contemporary method
frontier. Literature-backed candidates enter a later comparative study only after their
implementation provenance, licensing, mechanism checks, and tuning protocol have been
documented. A method is not elevated to a publication baseline merely because it conforms to
the same software API.

\subsection{Environment and scope}
The recorded run uses Python 3.12.3 and PyTorch 2.9.1 with ROCm 7.2 under WSL2 on an AMD
Radeon RX 7800 XT. Runtime is measured around the corresponding training operation. R001
includes P001 and separately labeled calibration evidence; none is treated as confirmatory.
The C001 design remains an unlocked local draft and may not begin until its decision rules,
manifests, and public prospective registration are finalized. Because the first
equivalence-margin calibration failed, C001 currently has no admissible equivalence claim.

The accelerator configuration is recorded because time and memory are properties of a
specific software--hardware stack, not universal method constants. PyTorch's ROCm backend
uses the “cuda” device namespace for compatibility even when CUDA itself is unavailable;
the environment record separately reports the ROCm runtime so this naming convention is not
misread as an NVIDIA execution claim. CPU execution remains supported for validation and
small tests, while future profiles may run unchanged on compatible CUDA, ROCm, Apple Silicon,
or cloud accelerators through the device abstraction.
